\documentclass[a4paper,11pt]{article}

\usepackage[ngerman]{babel}
%\usepackage[top=3cm,bottom=3cm,left=3cm,right=3cm]{geometry} %NICHT IN KEMIE2


\usepackage{microtype} %schöneres typesetting  mit [stretch=10] für nur 1% anstelle von 2%
\usepackage{lmodern} %Latin Modern FONT
\usepackage{pifont} %schönere Font

\usepackage{chemfig}
\usepackage[version=4]{mhchem}
\usepackage{amsmath}
\usepackage{nicefrac}
\usepackage{amssymb}
%\usepackage{mathtools}

\usepackage{titlesec} %Um Section, etc bearbeiten zu können
\usepackage{fancyhdr}
\usepackage{setspace}
\usepackage{enumitem} %enumerate
\usepackage{arydshln} %für dashline
\usepackage{multicol}
\usepackage{lipsum}
\usepackage[labelfont=bf]{caption}

\usepackage[many]{tcolorbox}

%\usepackage{pgfornament} %Scheene Ornamente
\usepackage{witharrows} %Arrows in align-Umgebung
%\usepackage{awesomebox} %Notizboxen, etc
%\usepackage{textpos} %Für zum Beispiel vertikale Texte am Rand


%\usepackage{wrapfig} %Text um Bilder
%\usepackage{varwidth} % Alternative für engere minipage


\usepackage{hyperref}


\tcbuselibrary{theorems}

%\titlespacing*{\section}{0pt}{10mm}{3mm}
%\titlespacing*{\subsection}{0pt}{6mm}{3mm}
%\titlespacing*{\subsubsection}{0pt}{8mm}{3mm}


\renewcommand{\ttdefault}{lmtt} %Schriftart wird geändert zu Latin Modern Typewriter


\hypersetup{hidelinks,
			colorlinks=true,
			linkcolor=black,
			citecolor=miudiutex,
			urlcolor=miugray2
			}

\setlength\fboxsep{0pt}				
\setlength\columnsep{20pt}
\setlength{\parindent}{0pt}
\setcounter{section}{0}






%\definecolor{miudiutex}{HTML}{2f3d39}
\definecolor{miudiutex}{HTML}{1d2926}
\definecolor{miugray}{HTML}{b4cccf}
%\definecolor{miugray2}{HTML}{4d7365}
\definecolor{miugray2}{HTML}{6F8C81}
\definecolor{beige}{HTML}{F4F8D3}
\definecolor{white2}{HTML}{f5f7f8}
\definecolor{red}{HTML}{cf1f5c}
\definecolor{red2}{HTML}{9C191B}
\definecolor{green}{HTML}{00A36C}
\definecolor{delphinidin2}{HTML}{315794}
\definecolor{delphinidin}{HTML}{8B9EBA}
\definecolor{uniblau}{HTML}{004e9f}
\definecolor{unigelb}{HTML}{fcba00}
\definecolor{unigrau}{HTML}{909085}
\definecolor{pelargonidin}{HTML}{b33807}
\definecolor{pelargonidin2}{HTML}{FFDC96}
\definecolor{cyanidin}{HTML}{ba0746}
\definecolor{paeonidin}{HTML}{d15ca6}
\definecolor{petunidin}{HTML}{B88495}
\definecolor{petunitext}{HTML}{963354}
\definecolor{malvidin}{HTML}{8a2d8c}
\definecolor{mytheorembg}{HTML}{F2F2F9}
\definecolor{mytheoremfr}{HTML}{00007B}
\definecolor{uniblau2}{HTML}{27548A}
\definecolor{unigelb2}{HTML}{DDA853}
\definecolor{unidunkelblau2}{HTML}{183B4E}
\definecolor{unibeige}{HTML}{F5EEDC}


\tikzset{>=Latex}

\raggedcolumns %wichtig für Multicolumns
\color{miudiutex}



\usetikzlibrary{chains,
                positioning,
                shapes.multipart,
                }


\setchemfig
{%	
	bond offset=2pt,
	atom sep=15pt, 
	cram width=3pt,
	cram dash width=0.8pt,
	cram dash sep=1.2pt,
	double bond sep=2.5pt, 
	bond style={line width=1.1pt,cap=round},
	baseline=(current bounding box.center),
	bond join=true
}
\setcharge{extra sep=3pt}





%================================
% SYMBOLS
%================================

\newcommand{\RA}{$\rightarrow$~}
\newcommand{\RL}{$\rightleftharpoons$~}
\newcommand{\HARP}{\ding{226}~}
%$\xrightarrow{2}$ %Pfeil mit 2 drüber


%================================
% SHORT COMMANDS (BOXES ; ETC)
%================================


\newcommand{\info}[1]{\begin{note}{\color{miudiutex}#1}\end{note}}
\newcommand{\qsbig}[1]{\begin{qstion}{}{}\begin{center}
			#1 \end{center}\end{qstion}}
\newcommand{\notiz}[1]{\begin{notizz}{}{}\begin{center}
			{\color{miugray2!50!miudiutex}#1} \end{center}\end{notizz}}
\newcommand{\klausur}[1]{\begin{klausurinfo}{}{}\begin{center}
			#1 \end{center}\end{klausurinfo}}
\newcommand{\dfn}[1]{\begin{defin}{}{}\begin{center}
			#1 \end{center}\end{defin}}
\newcommand{\gesetz}[2]{\begin{gestex}{}{}\begin{center}
			#1 \end{center}\end{gestex}}
\newcommand{\gergesetz}[1]{\begin{gesgex}{}{}\begin{center}
			#1 \end{center}\end{gesgex}}
\newcommand{\klaausur}[1]{\begin{klausurinfo}[width=.77\linewidth]{}{}\begin{center}
			\footnotesize #1 \end{center}\end{klausurinfo}}
\newcommand{\winfo}[1]{\begin{wichtig}{}{} \begin{center}
#1 \end{center} \vspace*{0mm} \end{wichtig}}
\newcommand{\zit}[2]{\begin{zitat}{#1}{}\small #2 \end{zitat}}



\newcommand{\unten}{\end{center} \tcblower \begin{center}}


\newcommand*\circled[1]{\tikz[baseline=(char.base)]{
            \node[shape=circle,draw,inner sep=0.5pt,miudiutex] (char) {\tiny #1};}}
            
\newcommand{\miusquare}{{\color{miugray2!50}\scriptsize $\blacksquare$}~}
\newcommand{\miusquarp}{{\color{petunidin!50}\scriptsize $\blacktriangleright$}~}
\newcommand{\niu}{\item[\miusquare]}
\newcommand{\nip}{\item[\miusquarp]}
\newcommand{\mrk}[1]{{\color{petunitext!55!red}\textsc{#1}}}

%%%%% ABSAETZE

\newcommand{\pps}{\par \vspace*{1mm}}
\newcommand{\pp}{\par \vspace*{3mm}}
\newcommand{\ppbig}{\par \vspace*{8mm}}
\newcommand{\pphuge}{\par \vspace*{10mm}}



%================================
% HEADER & FOOTER FUER UNI
%================================
\newcommand{\miusfoot}{
\fancypagestyle{plain}{%
  \fancyhf{}%
  \renewcommand{\footrulewidth}{0.0mm}%
  \fancyfoot[L]{\hspace*{-3cm} {\color{miugray2}\textbf{\thepage}}}%
  \renewcommand{\headrulewidth}{0mm}%
} 
\pagestyle{plain}
}

\newcommand{\miushead}[2]{
\miusquare #1 - 4.Semester - ELW \hfill 
\includegraphics[width=60pt]{{F://LaTeX-Files/shiet/unibo}}
\par \vspace*{-3mm}
\hrulefill
\par \vspace*{1mm}
\par \vspace*{-4mm}
\normalsize	
	\begin{flushright}
	\textbf{\miusquare #2}\par \vspace*{1.5mm}
	\small	
	\textit{Letzte Überarbeitung:} \par 
	\footnotesize \today
	\end{flushright}
	\normalsize
\par \vspace*{-8mm}
}

\newcommand{\sidebar}{
\AddToHook{shipout/background}{
  \begin{tikzpicture}[remember picture,overlay, shift={(current page.south west)}]
   \path[fill=miugray2!15](0,0) rectangle (3.5cm,\paperheight);
%   \path[fill=miugray2!8](0,0) rectangle (3.5cm,\paperheight);
   \path[draw,miudiutex!50] (3.5,0) to (3.5,\paperheight);
%   \path[draw,miudiutex!50,dashed] (3.5,0) to (3.5,\paperheight);  
%   \node[opacity=0.8](a)at(12,10){\includegraphics[width=9cm]{F://LaTeX-Files/pngs/girl01}};
  \end{tikzpicture} 
}
}


\newcommand{\oddsidebar}{
\AddToHook{shipout/background}{
   \put(0in,-\paperheight){%
      \ifodd\value{page}\relax
           \begin{tikzpicture}[remember picture,overlay, shift={(current page.south west)}]
   \path[fill=miugray2!15](0,0) rectangle (3.5cm,\paperheight);
      \path[draw,miudiutex!50] (3.5,0) to (3.5,\paperheight);
  \end{tikzpicture}
      \else
           \begin{tikzpicture}[remember picture,overlay, shift={(current page.south east)}]
   \path[fill=miugray2!15](-3,0) rectangle (3cm,\paperheight);
%      \path[draw,miudiutex!50] (-3.5,0) to (-3.5,\paperheight);
  \end{tikzpicture}
      \fi
   }
}
}
%================================
% ENVIRONMENTS
%================================

\newenvironment{ctab}[2]{%
				\begin{spacing}{#1}
					\begin{center}
					\begin{tabular}{#2}}
					{\end{tabular}
					\end{center}
				\end{spacing}
				}


%================================
% Farbiger-Background-Box
%================================



\newcommand{\farbig}[1]{
\begin{tcolorbox}[%
				hbox,
				colback=miugray2!25,
				boxsep=-2pt,
				boxrule=0pt,
				arc=0pt,
				nobeforeafter,
				tcbox raise base,
				shrink tight,
				extrude by=1pt
				]
				{\color{miudiutex}{#1}}
\end{tcolorbox}~}

%================================
% SIMPLE SCHWARZER RAHMEN BOX
%================================


\newcommand{\boxfr}[1]{
	\begin{tcolorbox}[
	colframe=miugray2,
	colback=white!95!miugray2,
	boxrule = 0.5pt, 
%	toprule=0pt, bottomrule=0pt,rightrule=0pt,leftrule=0pt,
	boxsep=0pt,
	arc=4pt,
	drop shadow southeast,
	enhanced
	]
	\color{miudiutex}
	\begin{center}
	#1
	\end{center}
	\end{tcolorbox}
}

%================================
% Antwort-Box
%================================

\newcommand{\antwort}[1]{\par 
	\begin{tcolorbox}[
	colback=gray!15,
	breakable,
	enhanced,
	frame hidden,
	arc=5pt,
	boxrule = 2pt,
%	borderline west = {2pt}{0pt}{petunidin},
%	drop shadow=miugray2!30
	]
	\color{miudiutex}	#1
	\end{tcolorbox}
}

%================================
% Frame-IT
%================================

\newcommand{\frameit}[1]{\par 
	\begin{tcolorbox}[
	colback=gray!10,
	breakable,
	enhanced,
	frame hidden,
	arc=5pt,
	boxrule = 0pt,
	borderline west = {4pt}{0pt}{petunidin},
%	drop shadow=miugray2!30
	]
		#1
	\end{tcolorbox}
}



%================================
% Question BOX
%================================

\makeatletter
\newtcbtheorem{qstion}{Frage}{enhanced,
	breakable,
	colback=white,
	colframe=delphinidin!100,
%	capture=hbox,
	boxsep=-3pt,
	attach boxed title to top left={yshift*=-\tcboxedtitleheight},
	fonttitle=\bfseries,
%	title={#2},
	boxed title size=title,
	boxed title style={%
			sharp corners,
			rounded corners=northwest,
			colback=tcbcolframe,
			boxrule=0pt,
		},
	underlay boxed title={%
			\path[fill=tcbcolframe] (title.south west)--(title.south east)
			to[out=0, in=180] ([xshift=5mm]title.east)--
			(title.center-|frame.east)
			[rounded corners=\kvtcb@arc] |-
			(frame.north) -| cycle;
		},
	#1
}{def}
\makeatother



%================================
% SIMPLE Question BOX
%================================

\newcommand{\qssmall}{
\begin{tcolorbox}[hbox,colback=miugray2!55,boxsep=-4pt,boxrule=0pt,arc=2pt,nobeforeafter, tcbox raise base,shrink tight,extrude by=3pt]
\textbf{\color{white}{?}}
\end{tcolorbox}~~}

\newcommand{\folie}[1]{
\begin{tcolorbox}[%
				hbox,
				colback=miugray2!20,
				boxsep=2pt,
				boxrule=0.5pt,
				arc=2pt,
				nobeforeafter,
				tcbox raise base,
				shrink tight,
				extrude by=3pt]
				{\color{miudiutex}{Folie #1}}
\end{tcolorbox}~}


%================================
% NOTIZ BOX
%================================

\makeatletter
\newtcbtheorem{notizz}{Notiz}{enhanced,
	breakable,
	coltitle=white,
	colback=petunidin!10,
	colframe=petunidin!50,
	attach boxed title to top left={yshift*=-\tcboxedtitleheight},
	fonttitle=\bfseries,
%	title={#2},
	boxed title size=title,
	boxed title style={%
			sharp corners,
			rounded corners=northwest,
			colback=tcbcolframe,
			boxrule=0pt,
		},
	underlay boxed title={%
			\path[fill=tcbcolframe] (title.south west)--(title.south east)
			to[out=0, in=180] ([xshift=5mm]title.east)--
			(title.center-|frame.east)
			[rounded corners=\kvtcb@arc] |-
			(frame.north) -| cycle;
		},
		#1
}{def}
\makeatother



%================================
% Question BOX: Klausur
%================================

\makeatletter
\newtcbtheorem[no counter]{klausurinfo}{Klausurrelevant}{enhanced,
%	breakable,
	colback=miugray2!10,
	colframe=miugray2!75,
	attach boxed title to top left={yshift*=-\tcboxedtitleheight},
	fonttitle=\bfseries,
	title={#2},
	arc=4pt,
	boxrule=1pt,
	boxed title size=title,
	boxed title style={%
			sharp corners,
			rounded corners=northwest,
			colback=tcbcolframe,
			boxrule=0pt,
			arc=4pt,
		},
	underlay boxed title={%
			\path[fill=tcbcolframe] (title.south west)--(title.south east)
			to[out=0, in=180] ([xshift=5mm]title.east)--
			(title.center-|frame.east)
			[rounded corners=\kvtcb@arc] |-
			(frame.north) -| cycle;
		},
	#1
}{def}
\makeatother


%================================
% DEFINITION
%================================

\makeatletter
\newtcbtheorem[no counter]{defin}{Definition}{enhanced,
	breakable,
	colback=gray!10,
	colframe=gray,
	attach boxed title to top left={yshift*=-\tcboxedtitleheight},
	fonttitle=\bfseries,
	title={#2},
	arc=4pt,
	boxrule=1pt,
	fontlower=\footnotesize,
	collower=miugray2!50!miudiutex,
	boxed title size=title,
	boxed title style={%
			sharp corners,
			rounded corners=northwest,
			colback=tcbcolframe,
			boxrule=0pt,
			arc=4pt,
		},
	underlay boxed title={%
			\path[fill=tcbcolframe] (title.south west)--(title.south east)
			to[out=0, in=180] ([xshift=5mm]title.east)--
			(title.center-|frame.east)
			[rounded corners=\kvtcb@arc] |-
			(frame.north) -| cycle;
		},
	#1
}{def}
\makeatother

%================================
% Gesetzes-Text EU
%================================

\makeatletter
\newtcbtheorem[no counter]{gestex}{EU - Verordnung}{enhanced,
	breakable,
	colback=gray!10,
	colframe={gray!65!delphinidin},
	attach boxed title to top left={yshift*=-\tcboxedtitleheight},
	fonttitle=\bfseries,
	arc=4pt,
	boxsep=10pt,
	boxrule=1pt,
	fontlower=\footnotesize,
	collower=miugray2!50!miudiutex,
	boxed title size=title,
	overlay={%
			\node[xshift=-1cm,yshift=-1mm] at (frame.north east)
			{\includegraphics[width=8mm]{F://LaTeX-Files/shiet/eu}};
%			\node[draw,miugray2] at (frame.north) {{#2}}; 
			},%
	boxed title style={%
			sharp corners,
			rounded corners=northwest,
			colback=tcbcolframe,
			boxrule=0pt,
			arc=4pt,
		},
	underlay boxed title={%
			\path[fill=tcbcolframe] (title.south west)--(title.south east)
			to[out=0, in=180] ([xshift=5mm]title.east)--
			(title.center-|frame.east)
			[rounded corners=\kvtcb@arc] |-
			(frame.north) -| cycle;
		},
	#1
}{def}
\makeatother


%================================
% Gesetzes-Text Ger
%================================

\makeatletter
\newtcbtheorem[no counter]{gesgex}{Gesetzestext}{enhanced,
	breakable,
	colback=gray!10,
	colframe={gray!95!pelargonidin},
	attach boxed title to top left={yshift*=-\tcboxedtitleheight},
	fonttitle=\bfseries,
	title={#2},
	arc=4pt,
	boxrule=1pt,
	fontlower=\footnotesize,
	collower=miugray2!50!miudiutex,
	boxed title size=title,
	overlay={\node[xshift=-1cm,yshift=-3mm] at (frame.north east)
			{\includegraphics[width=8mm]{F://LaTeX-Files/shiet/ger}}; 
			},
	boxed title style={%
			sharp corners,
			rounded corners=northwest,
			colback=tcbcolframe,
			boxrule=0pt,
			arc=4pt,
		},
	underlay boxed title={%
			\path[fill=tcbcolframe] (title.south west)--(title.south east)
			to[out=0, in=180] ([xshift=5mm]title.east)--
			(title.center-|frame.east)
			[rounded corners=\kvtcb@arc] |-
			(frame.north) -| cycle;
		},
	#1
}{def}
\makeatother


%================================
% Question BOX: Wichtig
%================================

%\makeatletter
%\newtcbtheorem[no counter]{wichtig}{Wichtige Info}{enhanced,
%	breakable,
%	colback=pelargonidin2!80,
%	colframe=red2!100,
%	fonttitle=\bfseries,
%	title={#2},
%	leftrule=5mm,
%	boxed title size=title,
%	boxed title style={%
%			sharp corners,
%			rounded corners=northwest,
%			colback=tcbcolframe,
%			boxrule=0pt,
%		},
%	underlay boxed title={%
%			\path[fill=tcbcolframe] (title.south west)--(title.south east)
%			to[out=0, in=180] ([yshift=1mm,xshift=7mm]title.east)--
%			([yshift=1mm]title.center-|frame.east)
%			[rounded corners=\kvtcb@arc] |-
%			(frame.north) -| cycle;
%		},
%		watermark tikz={\draw[line width=2mm] circle (1.1cm)
%		node{\fontfamily{ptm}\fontseries{b}\fontsize{20mm}{20mm}\selectfont !};}],
%	attach boxed title to top left={xshift=5mm,yshift*=-\tcboxedtitleheight},
%	watermark zoom=0.65,
%	  underlay={%
%    \path[draw=none] (interior.south west) rectangle node[white]{\Huge \textbf{!}} ([xshift=-4mm]interior.north west);
%    },
%	#1
%}{def}
%\makeatother


\makeatletter
\newtcbtheorem[no counter]{wichtig}{Wichtige Info}{enhanced,
	breakable,
	colback=pelargonidin2!30,
	colframe=petunidin!100,
	fonttitle=\footnotesize\bfseries,
	title={#2},
	leftrule=4mm,
	boxed title size=title,
	boxed title style={%
			sharp corners,
			rounded corners=northeast,
			colback=tcbcolframe,
			boxrule=0pt,
		},
	underlay boxed title={%
			\path[fill=tcbcolframe] (title.south east)--(title.south west)
			to[out=180, in=0] ([yshift=2mm,xshift=-7mm]title.west)--
			([yshift=2mm]title.center-|frame.west)
			[rounded corners=\kvtcb@arc] |-
			(frame.north) -| cycle;
		},
		watermark tikz={\draw[line width=2mm] circle (1.1cm)
		node{\fontfamily{ptm}\fontseries{b}\fontsize{20mm}{20mm}\selectfont !};},
	attach boxed title to top right={yshift*=-\tcboxedtitleheight},
	watermark zoom=1.45,
	watermark opacity=0.35,
	  underlay={%
    \path[draw=none] (interior.south west) rectangle node[white]{\LARGE \textbf{!}} ([xshift=-4.3mm]interior.north west);
    },
    #1
}{def}
\makeatother





%================================
% ZITAT BOX
%================================

\tcbuselibrary{theorems,skins,hooks}
\newtcbtheorem{zitat}{Zitat}
{%
	enhanced,
	breakable,
	colback = gray!10!white,
	frame hidden,
	boxrule = 0sp,
	borderline west = {2pt}{0pt}{petunidin},
	sharp corners,
	detach title,
	before upper = \tcbtitle\par\smallskip,
	coltitle = gray!50!black,
	fonttitle = \bfseries\sffamily,
	description font = \mdseries,
	separator sign none,
	segmentation style={solid, gray!50!black},
}
{th}



%\input{F://LaTeX-Files/shiet/vl.tex}
\newcommand{\metermilli}{\mskip3mu \text{mm}}
\newcommand{\cm}{\mskip3mu \text{cm}}
\newcommand{\m}{ \text{m}}
\newcommand{\lite}{ \text{l}}
\newcommand{\kg}{ \text{kg}}
\newcommand{\g}{ \text{g}}
\newcommand{\km}{ \text{km}}
\newcommand{\h}{ \text{h}}
\newcommand{\s}{ \text{s}}
\newcommand{\tonne}{ \text{t}}
\usepackage{pgfplots}
\usepackage[bottom=2cm,top=2cm,left=3cm,right=3cm]{geometry}
\title{PLMT Lösungen}

\begin{document}
\maketitle



\section*{Übung 01}
\subsection*{Aufgabe 01}
Ein Fluid durchströmt ein Rohr mit drei unterschiedlichen Querschnittsflächen. 
Im ersten Bereich beträgt die Querschnittsfläche 20 cm$^2$ und die Fließgeschwindigkeit des Fluids beträgt 3 cm/s.
\begin{itemize}
\item[a] Im zweiten Bereich beträgt die Querschnittsfläche 10 cm$^2$. Wie hoch ist die Geschwindigkeit des Fluids hier?
\item[b] Im dritten Bereich beträgt die Querschnittsfläche 25 cm$^2$. Wie hoch ist die Geschwindigkeit des Fluids hier?
\end{itemize}  

\ppbig
Kontinuitätsgleichung
\begin{align*}
A_1 \cdot v_1 = A_2 \cdot v_2 \\
v_2 = \frac{A_1 \cdot v_1}{A_2}
\end{align*}
Einsetzen a) und b):
\begin{align*}
a) &&v_2 &= \frac{20\cm^2 \cdot 3\mskip3mu \cm/\text{s}}{10\mskip3mu \cm^2} &&= 6\cm/\text{s} \\
b) &&v_2 &= \frac{20\cm^2 \cdot 3\mskip3mu \cm/\text{s}}{25\mskip3mu \cm^2} &&=2,4\cm/\text{s}
\end{align*}


\pagebreak
\subsection*{Aufgabe 02}
In einer Wasserstrahl-Vakuumpumpe, die wie ein Venturi-Rohr aufgebaut ist, fließen 4 L/min Wasser. Der Durchmesser der Zuleitung (und Ableitung) beträgt 14 mm, der Durchmesser an der engsten Stelle 3 mm. Der absolute Druck am Ausgang der Wasserstrahl-Vakuumpumpe entspricht einem Umgebungsdruck von 1 bar. Der geodätische Höhenunterschied kann vernachlässigt werden. 
Welcher Unterdruck wird an der engsten Stelle erreicht? Welcher absolute Druck herrscht an der engsten Stelle?
\ppbig
 

$\dot{V}=4 \mskip3mu \text{L/min} ; \qquad d_2 = 14 \metermilli ; \qquad d_1 = 3\metermilli ; \qquad \rho_{\text{H$_2$O}} = 1\g/\text{cm}^3$ 
\pp
Rechnung:
%\reversemarginpar
%\marginnote[Prinzipiell lässt sich $v_1$ auch über die selbe Methode wie $v_2$ errechnen. Ohne Fehler beim Runden sollte dann auch annähernd etwas mit 9,43 m/s heraus kommen...]{}[2.5cm]

\begin{align*}
\dot{V} &= v_2 \cdot A_2 \qquad \rightarrow v_2 = \frac{\dot{V}}{A_2} \\ \\
v_2 &= \frac{\frac{0,004 \mskip3mu \m^3}{60s}}{\pi \cdot (0,007~\m)^2} \\ 
%&\qquad \vdots \\
v_2 &= 0,43\mskip3mu \m/\text{s}\\ \\
v_1 &= v_2 \cdot \frac{A_2}{A_1}\\
&= v_2 \cdot \frac{(d_2)^2}{(d_1)^2} \\
&= 0,43\mskip3mu \text{m/s} \cdot \frac{14^2}{3^2} \\
&=9,36\mskip3mu \text{m/s} \\ \\
\end{align*}
Gesetz von Bernoulli:
\begin{equation}
p_1 + \frac{\rho \cdot v_1^2}{2} = p_2 + \frac{\rho \cdot v_2^2}{2}
\end{equation}
Daraus folgt:
\begin{align*}
p_2 - p_1 = \frac{\rho \cdot v_1^2}{2} -  \frac{\rho \cdot v_2^2}{2} 
\end{align*}
Werte einsetzen: 
\begin{align*}
p_2 - p_1 &= \frac{1000 \mskip3mu \text{kg/m}^3 \cdot (9,36\mskip3mu \text{m/s} )^2}{2} -  \frac{1000 \mskip3mu \text{kg/m}^3 \cdot (0,43 \mskip3mu \text{m/s} )^2}{2} \\
&= 43.712 \mskip3mu \frac{\text{kg}}{\m \cdot \s^2} \\
&= 43,7 \mskip3mu \text{kPa} \\ \\
p_1 &= p_2 - \Delta p \\
	&= (100 - 43,7)\text{kPa}\\ &= 56,3\mskip3mu \text{kPa}
\end{align*}
\pp
%\marginnote[Gesamtdruck = $p_0$]

Wobei $p_1$ der Druck an der engsten Stelle und $p_2$ der Druck am Ende des Rohres ist.\par
Absoluter Druck $p_1$ = 56,34 \text{kPa} \par 
Unterdruck $\Delta p$ = 43,7 \text{kPa} \par 




\pagebreak
\subsection*{Aufgabe 3}
Wie lange dauert es, bis Wasser aus einem Gefäß mit folgenden Eigenschaften vollständig abgeflossen ist?

\begin{itemize}
\item Situation A
\begin{itemize}
    \item  Querschnittsfläche des Gefäßes: 40 cm$^2$
    \item  Querschnittsfläche der Ausflussöffnung: 2 cm$^2$
    \item  Höhe des Flüssigkeitsspiegels über dem Ausfluss: 1 m
\end{itemize}
\item Situation B
\begin{itemize}
    \item  Querschnittsfläche des Gefäßes: 5 cm$^2$
    \item  Querschnittsfläche der Ausflussöffnung: 2 cm$^2$
    \item  Höhe des Flüssigkeitsspiegels über dem Ausfluss: 1 m
\end{itemize}
\item Situation C
\begin{itemize}
    \item  Querschnittsfläche des Gefäßes: 40 cm$^2$
    \item  Querschnittsfläche der Ausflussöffnung: 2 cm$^2$
    \item  Höhe des Flüssigkeitsspiegels über dem Ausfluss: 4 m
\end{itemize}
\item Situation D
\begin{itemize}
    \item  Querschnittsfläche des Gefäßes: 40 cm$^2$
    \item  Querschnittsfläche der Ausflussöffnung: 10 cm$^2$
    \item  Höhe des Flüssigkeitsspiegels über dem Ausfluss: 1 m
\end{itemize}
\end{itemize}

\ppbig
 

Torricelli:
\begin{equation}
t = \sqrt{\frac{2}{g}} \cdot \frac{A_1}{A_2} \cdot (\sqrt{h_a} - \sqrt{h_w})
\end{equation}
Für eine komplette Entleerung ($h_w$ = 0) :
\begin{align*}
t = \sqrt{\frac{2}{g}} \cdot \frac{A_1}{A_2} \cdot \sqrt{h_a}
\end{align*}
Einsetzen (Aufgabenteil a):
\begin{align*}
t &= 0,45 \frac{s}{\sqrt{m}}\cdot \frac{40cm^2}{2cm^2} \cdot \sqrt{1m} \\
 &= 9,03s 
\end{align*}
b) 1,13s ; c) 18,08s ; d) 1,81s \par 

\pagebreak
\subsection*{Aufgabe 4}
Wie viel Zeit braucht die Entleerung eines Rührkessels unter der Wirkung der \mbox{Erdschwere}, wenn der Bodenablass die NW 80mm aufweist? Der Rührkessel ist zu Beginn mit \mbox{10.000 L} Wasser gefüllt. Der Innendurchmesser des Rührkessels beträgt 2,4 m. Der gewölbte Boden des Klöppers ist bei der Berechnung zu vernachlässigen, ebenso der Druckverlust beim Einströmen in die Rohrleitung.
\ppbig
 


\begin{align*}
&A_1 &&= \pi \cdot (\frac{2,4\mskip3mu \m}{2})^2 &&= 4,52\mskip3mu \m^2 \\
&A_2 &&= \pi \cdot (\frac{0,08\mskip3mu \m}{2})^2 &&= 0,005\mskip3mu \m^2 \\
&h_a &&= \frac{V}{A_1} = \frac{10\mskip3mu \m^3}{4,52m^2} &&= 2,21\mskip3mu \m^2 
\end{align*}

\ppbig
Torricelli:
\begin{equation}
t = \sqrt{\frac{2}{g}} \cdot \frac{A_1}{A_2} \cdot (\sqrt{h_a} - \sqrt{h_w})
\end{equation}

Einsetzen:
\begin{align*}
t &= \sqrt{\frac{2}{g}} \cdot \frac{4,52\mskip3mu \m^2}{0,005\mskip3mu \m^2} \cdot \sqrt{2,21\mskip3mu \m^2} \\
&= 604s
\end{align*}







\pagebreak
\subsection*{Aufgabe 5}
Berechnen Sie die Reynoldszahl für Wasser bei 0°C. Ist die Strömung jeweils laminar oder turbulent?
\begin{itemize}
       \item[a.] Geschwindigkeit 5 cm/s, Rohrinnendurchmesser 30 cm
       \item[b.] Geschwindigkeit 1 cm/s, Rohrinnendurchmesser 30 cm
       \item[c.] Geschwindigkeit 5 cm/s, Rohrinnendurchmesser 10 cm
\end{itemize}

\ppbig

 

Kritische Reynoldszahl für turbulente oder laminare Strömung:
\begin{equation}
\text{Re}_{\text{krit}} = \frac{\rho \cdot d \cdot v}{\eta} = 2300
\end{equation}
\reversemarginpar
%\marginnote[Vorsicht! Pascal ist in SI-Einheiten N/m$^2$ \par \RA Rest auch unbedingt in SI umrechnen]

Einsetzen:
\begin{align*}
&\text{a)} && \frac{1.000\mskip3mu \kg/\m^3 \cdot 0,3\mskip3mu \m \cdot 0,05\mskip3mu \m/\s}{1.790 \cdot 10^{-6}\mskip3mu \text{P} \cdot \s} &&\approx 8379,9 \\
&\text{b)} && \frac{1.000\mskip3mu \kg/\m^3 \cdot 0,3\mskip3mu \m \cdot 0,01\mskip3mu \m/\s}{1.790 \cdot 10^{-6}\mskip3mu \text{P} \cdot \s} &&\approx 1676 \\
&\text{c)} && \frac{1.000\mskip3mu \kg/\m^3 \cdot 0,1\mskip3mu \m \cdot 0,05\mskip3mu \m/\s}{1.790 \cdot 10^{-6}\mskip3mu \text{P} \cdot \s} &&\approx 2793,3 \\
\end{align*}

Oberhalb der kritischen Reynolds-Zahl 2.300 erfolgt der Übergang zu turbulenten Strömungen


\pagebreak





\section*{Übung 02}
\subsection*{Aufgabe 01}
Öl der Viskosität 0,031 kg/m*s und der Dichte 850 kg/m$^3$ soll mit 13 m$^3$/h bei einem
zulässigen Druckabfall von 0,1 bar laminar durch ein kreisrundes Rohr der Länge
l = 100m geleitet werden. \par 
Wie groß ist die Nennweite? Ist die Annahme einer laminaren Strömungsform
berechtigt?

Umformen:
\begin{align*}
P &= \frac{8 \eta \cdot \Delta l}{\pi \cdot R^4} \cdot \dot{V}^2 \\
\dot{V} \cdot \Delta p &=  \frac{8 \eta \cdot \Delta l}{\pi \cdot R^4} \cdot \dot{V}^2 \\
 \Delta p &=  \frac{8 \eta \cdot \Delta l}{\pi \cdot R^4} \cdot \dot{V} \\
R^4 &=  \frac{8 \eta \cdot \Delta l}{\pi \cdot  \Delta p} \cdot \dot{V} \\
R &= \sqrt[4]{\frac{8 \eta \cdot \Delta l}{\pi \cdot  \Delta p} \cdot \dot{V}}
\end{align*}
Werte einsetzen:
\begin{align*}
R &= \sqrt[4]{\frac{8 \cdot 0,031 ~\text{kg/ms} \cdot 100~\text{m}}{\pi \cdot  0,1 \text{bar}} \cdot 13~\text{m}^3/\text{h}} \\
R &= \sqrt[4]{\frac{8 \cdot 0,031 ~\text{kg/ms} \cdot 100~\text{m}}{\pi \cdot  10.000 ~\text{kg/ms}^2} \cdot 0,00361~\text{m}^3/\text{s}}\\
R &\approx 0,041~\text{m} \\
d &\approx 0,082~\text{m} \\
d &\approx 8,2~\text{cm}
\end{align*}
Schauen ob Strömung laminar oder turbulent:
\begin{align*}
\text{Re}_{\text{krit}} = \frac{\rho d v}{\eta} = 2300
\end{align*}
$v$ ausrechnen:
\begin{align*}
v = \frac{\dot{V}}{A} \rightarrow v &= \frac{13~\text{m}^3/\text{h}}{\pi \cdot (0,041~\text{m})^2} \\
&\approx 2461,65~\text{m/h} \\
&\approx 0,684~\text{m/s} 
\end{align*}
Werte einsetzen:
\begin{align*}
\frac{850~\text{kg/m}^3 \cdot 0,082~\text{m} \cdot 0,684~\text{m/s}}{0,031~\text{kg/ms}} \approx 1537,9 \rightarrow {laminar}
\end{align*}



\pagebreak

\subsection*{Aufgabe 02}
Wasser von 20 °C fließt mit einem Volumenstrom von 600 m$^3$/h durch ein Stahlrohr mit der Nennweite 200 mm. \par 
Wie groß ist der Druckverlust bei einer Leitungslänge von 1000 m?

$\xi$ von Stahlrohr $\approx$ 0,0113 \par 
$\overline{v}$ ausrechnen:
\begin{align*}
\overline{v} = \frac{\dot{V}}{A} \rightarrow \overline{v} &= \frac{600~ \m^3/h}{\pi \cdot (0,1~\text{m})^2} \\
 &\approx 19.098,6~\text{m/h} \\
 &\approx 5,3~\text{m/s}  \\ \\
\end{align*}
Wir nutzen die Formel für Druckabfall im Rohr:
\begin{align*}
\Delta p &= \xi \cdot \frac{l}{d} \cdot \frac{1}{2} \rho \cdot \overline{v}^2\\
\end{align*}
Werte einsetzen:
\begin{align*}
\Delta p &= 0,0113 \cdot \frac{1000~\text{m}}{0,2~\text{m}} \cdot \frac{1}{2} \cdot 998,20 ~\text{kg}/~\text{m}^3 \cdot (5,3~m/s)^2 &&\approx 792.114~\text{Pa}\\
& &&\approx 7,92 \text{bar}
\end{align*}


\pagebreak

\subsection*{Aufgabe 03}
Der Volumenstrom von Bier, das durch ein Rohr fließt, beträgt 1,8 L/s. Der
Innendurchmesser des Rohres beträgt 3 cm. Das Bier hat eine Dichte von 1100 kg/m$^3$.
Berechnen Sie die durchschnittliche Geschwindigkeit von Bier in m/s und seinen Massendurchfluss in kg/s. \par Wie hoch ist die Geschwindigkeit bei gleichem Volumenstrom, wenn ein anderes Rohr mit einem Durchmesser von 1,5 cm verwendet wird?

$\dot{V} = 1,8 ~$L$/$s $ = 0,0018~$m$^3/$s \par 
$r=0,015~$m \pp
Massendurchfluss:
\begin{align*}
0,0018~\m^3/\s \cdot 1100~\kg/\m^3 = 1,98~\kg/\s
\end{align*}
Mittlere Geschwindigkeiten für versch. Durchmesser:
\begin{align*}
\overline{v} = \frac{\dot{V}}{A} \rightarrow \qquad \overline{v} &= \frac{0,0018~\m^3/s}{\pi \cdot (0,015~\text{m})^2} \qquad = 2,546~\m/\s \\
\overline{v} &= \frac{0,0018~\m^3/\s}{\pi \cdot (0,0075~\text{m})^2} \qquad = 10,19~\m/\s \\
\end{align*}



\subsection*{Aufgabe 04}
Ein Edelstahltank mit einem Durchmesser von 3 m enthält Wein. Der Tank wird bis zu
einer Höhe von 5 m befüllt. Eine Auslassöffnung mit einem Durchmesser von 10 cm wird
geöffnet, um den Wein abzulassen. \par Berechnen Sie die Austragsgeschwindigkeit des
Weins, unter der Annahme, dass der Fluss gleichmäßig und reibungslos ist. Wie lange
dauert es, bis der Tank vollständig entleert ist?


Toricelli:
\begin{align*}
t = \sqrt{\frac{2}{g}} \cdot \frac{A_1}{A_2} \cdot (\sqrt{h_{\alpha}}-\sqrt{h_{\omega}})
\end{align*}
Werte einsetzen:
\begin{align*}
t = \sqrt{\frac{2}{9,81~\m/\s^2}} \cdot \frac{\pi \cdot (1,5~\m)^2}{\pi \cdot (0,05~\m)^2} \cdot \sqrt{5~\m} \qquad = 908,67~\s
\end{align*}



\pagebreak
\subsection*{Aufgabe 05}
Zeichnen Sie die Fließ- und Viskositätskurve einer newton'schen Flüssigkeit, einer
strukturviskosen Flüssigkeit und einer dilatanten Flüssigkeit.
\ppbig

\begin{tikzpicture}
\node[petunitext,scale=1.1,align=left](12)at(0,2){Newtonsch};
\node[scale=0.7](a) at(0,0){
\begin{tikzpicture}
\begin{axis}[
	title=Fließkurve,
	ylabel=$\tau$,
	xlabel=$D$,
	xmin=0,
	ymin=0,
	xmax=3.5,
	ymax=3.5
	]
\addplot[mark=none,miugray2] {x};
\end{axis}
\end{tikzpicture}
}; \hspace*{4cm}
\node[scale=0.7](b)at(4,0){
\begin{tikzpicture}

\begin{axis}[
	title= Viskositätskurve,
	ylabel=$\eta$,
	xlabel=$D$,
	xmin=0,
	ymin=0,
	xmax=3.5,
	ymax=3.5
	]
\addplot[mark=none,miugray2] {2};
\end{axis}
\end{tikzpicture}
};
\end{tikzpicture}
\pp
\ppbig
\begin{tikzpicture}
\node[petunitext,scale=1.1,align=left](12)at(0,2){Strukturviskos};
\node[scale=0.7](a) at(0,0){
\begin{tikzpicture}
\begin{axis}[
	title=Fließkurve,
	ylabel=$\tau$,
	xlabel=$D$,
	xmin=0.2,
	ymin=0,
	xmax=2,
	ymax=3.5
	]
\addplot[mark=none,miugray2] {-(x-2)^2+3.5};
\end{axis}
\end{tikzpicture}
}; \hspace*{4cm}
\node[scale=0.7](b)at(4,0){
\begin{tikzpicture}

\begin{axis}[
	title= Viskositätskurve,
	xlabel=$D$,
	ylabel=$\eta$,
	xmin=0.7,
	ymin=0,
	xmax=3.5,
	ymax=3.5
	]
\addplot[mark=none,miugray2,domain=0.1:3.5] {1/(x^2)+1};\end{axis}
\end{tikzpicture}
};
\end{tikzpicture}
\pp
\ppbig
\begin{tikzpicture}
\node[petunitext,scale=1.1,align=left](12)at(0,2){Dilatant};
\node[scale=0.7](a) at(0,0){
\begin{tikzpicture}
\begin{axis}[
	title=Fließkurve,
	ylabel=$\tau$,
	xlabel=$D$,
	xmin=0,
	ymin=0,
	xmax=6,
	ymax=6
	]
\addplot[mark=none,miugray2] {0.25*x^2};
\end{axis}
\end{tikzpicture}
}; \hspace*{4cm}
\node[scale=0.7](b)at(4,0){
\begin{tikzpicture}

\begin{axis}[
	title= Viskositätskurve,
	ylabel=$\eta$,
	xlabel=$D$,
	xmin=0,
	ymin=0,
	xmax=6,
	ymax=6
	]
\addplot[mark=none,miugray2] {0.25*x^2+2};
\end{axis}
\end{tikzpicture}
};
\end{tikzpicture}








	
\end{document}









\end{document}